% vim: set foldmethod=marker foldlevel=0:

\documentclass[a4paper]{article}
\usepackage[UKenglish]{babel}

\usepackage{preamble}

% \renewcommand{\thesubsection}{Q\arabic{section}~(\roman{subsection})}

\fancyhead[L]{MA2K4 Assignment 4}
\title{MA2K4 Numerical Methods and Computing, Assignment 4}
\colorlet{questionbodycolor}{green!50}

\begin{document}

\maketitle

\setlength{\parindent}{0em}
\setlength{\parskip}{1em}

% {{{ Q1
\question{1}

\begin{questionbody}
Let $f_n = f(t_n, u_n)$ and consider the one-step method \[
u_{n+1} = u_n + \alpha h f_n + \beta h f(t_n + \gamma h, u_n + \gamma h f_n),
\] where $\alpha, \beta, \gamma \in \R$, and $h > 0$ is the stepsize.
\begin{enumerate}[(a)]
\item Compute the local truncation error of the method in terms of $\alpha$, $\beta$, and $\gamma$. From this, show that this method is consistent if and only if $\alpha + \beta = 1$.

\item With the same computation to one order further, show that the method is of second order if $\beta \gamma = \f12$.

\item Show that the order of the method cannot be higher than 2.
\end{enumerate}
\end{questionbody}

\subsection{~} % 1.a

Answer

\subsection{~} % 1.b

Answer

\subsection{~} % 1.c

Answer

% }}}

% {{{ Q2
\newquestion{2}

\begin{questionbody}
Derive the region of absolute stability for the Heun method \[
u_{n+1} = u_n + \f12 h \big( f_n + f(t_n + h, u_n + h f_n) \big)
\] by stating an inequality which needs to be fulfilled whenever the method is absolutely stable.
\end{questionbody}

Answer

% }}}

% {{{ Q3
\newquestion{3}

\begin{questionbody}
The $\theta$-method is given by \[
u_{n+1} = u_n + h \big( (1 - \theta) f_n + \theta f_{n+1} \big)
\] for $\theta \in [0, 1]$. Show that the $\theta$-method is $A$-stable if and only if $\theta \ge \f12$.
\end{questionbody}

Answer

% }}}

% {{{ Q4
\newquestion{4}

\begin{questionbody}
In this problem, we will discuss different Runge--Kutta methods in light of the order and consistency conditions derived in the lecture. % chktex 8
\begin{enumerate}[(a)]
\item Show that the Runge--Kutta method with Butcher tableau \[ % chktex 8
\begin{array}{c|ccc} % chktex 44
0 \\
1/2 & 1/2 \\
3/4 & & 3/4 \\
\hline % chktex 44
& 2/9 & 1/3 & 4/9
\end{array}
\] has maximal order.

\item Construct a 1-stage Runge--Kutta method of highest possible order and write down its Butcher tableau. Is the answer unique? % chktex 8

\item Construct a general second order explicit 2-stage Runge--Kutta method and write down its Butcher tableau. Express everything in the minimal number of parameters, so that the method fulfils all required conditions. % chktex 8
\end{enumerate}
\end{questionbody}

\subsection{~} % 4.a

Answer

\subsection{~} % 4.b

Answer

\subsection{~} % 4.c

Answer

% }}}

% {{{ Q5
\newquestion{5}

\begin{questionbody}
For a general 2-stage implicit Runge--Kutta method applied to the test problem $y'(t) = \lambda y(t)$, show that % chktex 8
\begin{align*}
K_1 &= \big( 1 + \lambda h (a_{12} - a_{22}) \big) \lambda u_n / \Delta, \\
K_2 &= \big( 1 + \lambda h (a_{21} - a_{11}) \big) \lambda u_n / \Delta,
\end{align*}
where $\Delta$ is the determinant of the matrix $I - \lambda h A$, where $A$ is taken from the Butcher tableau of the method.
\end{questionbody}

Answer

% }}}

% {{{ Q6
\newquestion{6}

\begin{questionbody}
Consider the 2-stage implicit Runge--Kutta method given by the Butcher tableau \[ % chktex 8
\begin{array}{c|cc} % chktex 44
(3 - \sqrt 3) / 6 & 1/4 & (3 - 2 \sqrt 3) / 12 \\
(3 + \sqrt 3) / 6 & (3 + 2 \sqrt 3) / 12 & 1/4 \\
\hline % chktex 44
& 1/2 & 1/2
\end{array}
\] Deduce that, when applied to the test problem, we have $u_{n+1} = R (\lambda h) u_n$, where \[
R (\lambda h) = \f{1 + \f12 \lambda h + \f1{12} \lambda^2 h^2}{1 - \f12 \lambda h + \f1{12} \lambda^2 h^2}.
\]
\end{questionbody}

Answer

% }}}

\end{document}
